\documentclass[10pt]{article}
\usepackage{tikz}
\usetikzlibrary{arrows}
\usepackage{braket}
\usepackage{algorithm}
\usepackage{algpseudocode}
\usepackage{soul}
\usepackage{float}
\usepackage{tabularx}
\usepackage{mathtools}
\usepackage{amsmath, amssymb, amsthm}
\usepackage{enumitem}
\usepackage{geometry}
\usepackage{fancyhdr}
\usepackage{titlesec}
\usepackage{xltabular}
\usepackage{listings, xcolor}
\usepackage{courier}
\usepackage{color}
\usepackage{hyperref}
\geometry{margin=1in}

\pagestyle{fancy}
\fancyhf{}
\cfoot{\thepage}
\title{SFWRENG 3S03 - Assignment 2}
\author{Matthew Farah, \href{mailto:farahm11@mcmaster.ca}{$\langle$farahm11@mcmaster.ca$\rangle$}, 400468251}
\date{\today}

\hypersetup{
	colorlinks=true,
	linkcolor=blue,
	filecolor=magenta,
	urlcolor=blue,
	citecolor=blue,
	anchorcolor=blue
}

\lstset{
	tabsize = 4,
	showstringspaces = false,
	numbers = left,
	commentstyle = \color{teal},
	keywordstyle = \color{blue},
	stringstyle = \color{red},
	rulecolor = \color{black},
	basicstyle = \ttfamily,
	breaklines = true,
	numberstyle = \small,
}

\fancyhead{}
\fancyhead[L]{Matthew Farah, \href{mailto:farahm11@mcmaster.ca}{$\langle$farahm11@mcmaster.ca$\rangle$}, 400468251}
\fancyhead[R]{SFWRENG 3S03 - Assignment 2}

\AtBeginEnvironment{align}{\setcounter{equation}{0}}

\setlength{\parindent}{0cm}

\setcounter{tocdepth}{3}
\hypersetup{bookmarksopen=true, bookmarksopenlevel=2, bookmarksdepth=3}

\newcounter{chapter}
\renewcommand{\thechapter}{\Roman{chapter}}

\newcounter{question}
\renewcommand{\thequestion}{\arabic{question}}

\newenvironment{question}[2]{
	\refstepcounter{question}
	\newpage
	\phantomsection
	\addcontentsline{toc}{section}{\protect{\large{Question \thequestion \, - #1}}}
	\vspace{1em}
	\par
	\large{[#2 marks]}\par
	\large\textbf{Question \thequestion}
	\vspace{.2cm}
	\par
	\setcounter{subquestion}{0}
	\setcounter{step}{0}
}{\vspace{.5em}}

\newenvironment{solution}{
	\vspace{1em}\par\large\textbf{{Solution:}}\par\vspace{1em}
}{
	\vspace{1em}
}

\newcounter{subquestion}
\renewcommand{\thesubquestion}{\alph{subquestion}}

\newenvironment{subquestion}{
	\refstepcounter{subquestion}
	\phantomsection
	\vspace{1em}
	\par\large\textbf{(\thesubquestion)}\hspace{-1em}
	\setcounter{step}{0}
	\addcontentsline{toc}{subsection}{\protect{\large{Part \thequestion.\thesubquestion}}}
}{
	\par
	\vspace{1em}
}

\makeatother

\makeatletter

\newcounter{step}
\renewcommand{\thestep}{\Roman{step}}

\newenvironment{step}{
	\refstepcounter{step}
	\vspace{1.5em}\par\noindent\large\textbf{(\thestep)}
}{
	\par
	\vspace{1.5em}
}

\makeatother

\newcolumntype{Y}{>{\centering\arraybackslash}X}
\renewcommand{\arraystretch}{1.8}

\fontsize{10pt}{14pt}\selectfont

\begin{document}

\maketitle

\tableofcontents
\newpage

\begin{question}{Basic deployment pipeline – Modeling with Petri Nets}{3.5}
	Incoming update requests are submitted to the deployment service. Each request
	corresponds to a deployment batch, which must pass through a sequence of
	processing stages before being installed on devices. \\

	The pipeline includes the following shared resources:
	\begin{itemize}
		\item Validation Service – verifies update packages before deployment.
		\item Deployment Manager – coordinates sending updates to devices.
		\item Two Deployment Workers – workers that push updates to devices in parallel.
		\item Staging buffer – can hold at most three deployment batches simultaneously. \\
	\end{itemize}

	The validation service, the deployment manager, and the deployment workers
	process only one batch at a time. \\

	Each deployment batch goes through the following stages:
	\begin{itemize}
		\item Submission – an update request enters the system.
		\item Staging – the batch waits in the staging buffer.
		\item Validation – the validation service checks the update package.
		\item Dispatch – the deployment manager assigns the batch to an available worker.
		\item Deployment – a worker performs the update rollout.
		\item Completion – the batch finishes deployment successfully. \\
	\end{itemize}

	Model the system in a Petri net.
\end{question}

\begin{solution}

\end{solution}

\begin{question}{Safety and other useful properties – Formalizing requirements}{3.5}
	Formulate CTL queries to capture the following requirements. For each requirement:
	\begin{enumerate}
		\item Formulate the CTL expression and explain it in your report;
		\item Implement it in TAPAAL;
		\item Disclose the result of the query in your report and interpret it.
	\end{enumerate}
\end{question}

\begin{subquestion}
	Buffer safety – The system never exceeds the staging buffer capacity.
\end{subquestion}

\begin{solution}

\end{solution}

\begin{subquestion}
	Resource exclusivity – The validation service and deployment manager are
	never used by more than one batch at the same time.
\end{subquestion}

\begin{solution}

\end{solution}

\begin{subquestion}
	Deadlock possibility – Determine whether the system can reach a marking
	where no transitions are enabled.
\end{subquestion}

\begin{solution}

\end{solution}

\begin{subquestion}
	Pipeline progress – If a batch enters the system, is it always possible for it to
	eventually complete deployment?
\end{subquestion}

\begin{solution}

\end{solution}

\begin{question}{Critical updates – Working with priorities}{3}
	In practice, some updates must be deployed with higher urgency than regular feature
	updates. Pertinent examples include security and safety patches. \\

	Extend your model to be able to model two types of deployment batches: regular and
	critical updates. The system must enforce the following operational policy: whenever
	at least one critical update batch is waiting to be dispatched, the deployment
	manager must assign it to a worker before any regular update batch. \\

	Implement this policy in your Petri net using inhibitor arcs. (The
	current version of TAPAAL does not support priority transitions.) \\

	Disclose your extended model in your report and identify the key changes.
\end{question}

\begin{solution}

\end{solution}

\begin{question}{Formalizing Requirements}{3}
	Formulate CTL queries to capture the following requirements. For each requirement:
	\begin{enumerate}
		\item Formulate the CTL expression and explain it in your report;
		\item Implement it in TAPAAL;
		\item Disclose the result of the query in your report and interpret it.
	\end{enumerate}
\end{question}

\begin{subquestion}
	Priority enforcement – When critical updates are waiting for deployment,
	regular updates cannot begin deployment.
\end{subquestion}

\begin{solution}

\end{solution}

\begin{subquestion}
	Critical update progress – Critical updates that enter the system eventually
	complete deployment.
\end{subquestion}

\begin{solution}

\end{solution}

\begin{subquestion}
	Regular update starvation – Determine whether regular updates may be
	indefinitely delayed if critical updates keep arriving.
\end{subquestion}

\begin{solution}

\end{solution}

\begin{question}{Reflection}{5}
	Reflect on this assignment and the overall module (MBT and Petri nets) briefly, e.g., by
	elaborating on the following questions, or your own points. \\

	\begin{itemize}
		\item How do you see verification being useful in your future industry job? What did
		you learn about the complexity of proving correctness (verification) vs
		increasing our confidence in the correctness (validation) of the program?
		Which one would you use in your daily professional activities and why?
		\item What is the role of a tool in mathematically heavy methods, such as verification
		with Petri nets? Does the existence of a good tool make a difference in the
		adoption of the formal verification techniques, or is it mostly irrelevant?
		\item What is the link between verification and validation.
		\item What was the hardest part of this assignment and why?
	\end{itemize}
\end{question}

\begin{solution}

\end{solution}

\end{document}
